\documentclass[11pt,a4paper]{article}
%\usepackage[utf8]{inputenc}
%\usepackage[T1]{fontenc}
\usepackage{fontspec}
\usepackage[french]{babel}
\usepackage{amsmath,amssymb}
\usepackage{geometry}
\usepackage{fancyhdr}
\setlength{\headheight}{15pt}
\usepackage{enumitem}
\usepackage{multicol}
\usepackage{tikz}
\usepackage{pgfplots}
\pgfplotsset{compat=1.18}
\usepackage{xcolor}
\usepackage{graphicx}
\usepackage{array}
\usepackage{fontawesome5}
\usepackage{tabularx}
\usepackage[most]{tcolorbox}

% Définition des couleurs modernes
\definecolor{primary}{RGB}{41, 128, 185}      % Bleu moderne
\definecolor{secondary}{RGB}{52, 73, 94}      % Gris foncé
\definecolor{accent}{RGB}{46, 204, 113}       % Vert menthe
\definecolor{warning}{RGB}{231, 76, 60}       % Rouge moderne
\definecolor{light}{RGB}{236, 240, 241}       % Gris très clair
\definecolor{bglight}{RGB}{247, 249, 250}     % Fond clair

\DeclareRobustCommand{\rep}[1]{\begingroup\color{warning}\ifmmode #1\else \textbf{#1}\fi\endgroup}

% Configuration de la page
\geometry{
	top=1.5cm,
	bottom=2cm,
	left=1.8cm,
	right=1.8cm
}

% Configuration tcolorbox
\tcbuselibrary{skins,breakable}

% Boîte pour les exercices
\newtcolorbox{exercisebox}[2][]{
	colback=white,
	colframe=primary,
	fonttitle=\bfseries\large,
	title={#2},
	arc=2mm,
	boxrule=1.5pt,
	left=15pt,
	right=15pt,
	top=12pt,
	bottom=12pt,
	boxsep=0pt,
	breakable,
	enhanced,
	width=\linewidth,
	grow to left by=3mm,
	grow to right by=3mm,
	#1
}

% Boîte pour les consignes
\newtcolorbox{consignebox}{
	colback=light!30,
	colframe=secondary,
	arc=3mm,
	boxrule=0pt,
	leftrule=4pt,
	left=15pt,
	right=15pt,
	top=12pt,
	bottom=12pt,
	enhanced,
	width=\linewidth,
	grow to left by=3mm,
	grow to right by=3mm
}

% En-tête et pied de page personnalisés
\pagestyle{fancy}
\fancyhf{}
\fancyhead[L]{\color{secondary}\textbf{6\textsuperscript{e}} | Mathématiques}
\fancyhead[R]{\color{secondary}\textbf{Nombres Entiers}}
\fancyfoot[C]{\color{secondary}\thepage}
\renewcommand{\headrulewidth}{0pt}
\renewcommand{\footrulewidth}{0pt}

% Commandes personnalisées
\newcommand{\points}[1]{\hfill \colorbox{primary}{\color{white}\textbf{\,#1 pts\,}}}
\newcommand{\badge}[2]{\tikz[baseline=(char.base)]{\node[shape=rectangle,draw,thick,rounded corners=2mm,inner sep=2pt,fill=#1,text=white,font=\footnotesize\bfseries] (char) {#2};}}

% Pointillés et zones de réponse
\newcommand{\dotline}[1]{\makebox[#1][s]{\dotfill}}
\newcolumntype{L}[1]{>{\hsize=#1\hsize\arraybackslash}X}

\newcommand{\zonereponse}[1]{%
	\par\noindent
	\colorbox{light}{%
		\parbox{\dimexpr\linewidth-2\fboxsep\relax}{%
			\vspace{#1}%
		}%
	}%
	\par
}

\newcommand{\zonereponseLignes}[2][0.9cm]{%
	\par\noindent
	\colorbox{light}{%
		\parbox{\dimexpr\linewidth-2\fboxsep\relax}{%
			\setlength{\parskip}{#1}%
			\foreach \i in {1,...,#2}{\noindent\dotfill\par}%
		}%
	}\par
}

\begin{document}
	
	% En-tête moderne avec TikZ
	\begin{tikzpicture}[remember picture,overlay]
		\fill[primary] (current page.north west) rectangle ([yshift=-3.5cm]current page.north east);
		\node[anchor=north west,text=white,font=\Huge\bfseries] at ([xshift=2cm,yshift=-1cm]current page.north west) {Contrôle \no 1};
		\node[anchor=north west,text=white,font=\Large] at ([xshift=2cm,yshift=-2cm]current page.north west) {Les Nombres Entiers};
		\node[anchor=north east,text=white,font=\large] at ([xshift=-2cm,yshift=-1.5cm]current page.north east) {\faIcon{clock} 50 minutes};
		\node[anchor=north east,text=white,font=\large] at ([xshift=-2cm,yshift=-2.3cm]current page.north east) {\faIcon{calculator} Calculatrice interdite};
	\end{tikzpicture}
	
	\vspace{1.5cm}
	
	% Informations élève
	\begin{tcolorbox}[
		colback=bglight,colframe=secondary,arc=3mm,boxrule=1pt,
		left=15pt,right=15pt,top=10pt,bottom=10pt,
		enhanced,
		width=\linewidth,
		grow to left by=5mm,
		grow to right by=5mm
		]
		\setlength{\tabcolsep}{4pt}
		\begin{tabularx}{\linewidth}{@{}lX lX l@{}}
			\multicolumn{2}{@{}l}{\faIcon{user} \textbf{NOM :} \dotline{6cm}} &
			\multicolumn{2}{l@{}}{\textbf{Prénom :} \dotline{3cm}} &
			\\[0.5cm]
			\textbf{Classe :} & \dotfill &
			\faIcon{calendar} \textbf{Date :} & \dotfill &
			\fcolorbox{secondary}{white}{\textbf{NOTE : }\makebox[1.4cm]{\rule{0pt}{1.1em}}\textbf{ / 20}}
		\end{tabularx}
	\end{tcolorbox}
	
	\vspace{0.2cm}
	
	% Barème
	\begin{center}
		\badge{primary}{Ex.1: 4pts} \quad
		\badge{primary}{Ex.2: 5pts} \quad
		\badge{primary}{Ex.3: 4pts} \quad
		\badge{primary}{Ex.4: 4pts} \quad
		\badge{primary}{Ex.5: 3pts} \quad
		\badge{accent}{Bonus: +2pts}
	\end{center}
	
	\vspace{0.2cm}
	
	% Consignes
	\begin{consignebox}
		\faIcon{info-circle} \textbf{CONSIGNES IMPORTANTES}
		\begin{itemize}[leftmargin=*,noitemsep,topsep=5pt]
			\item[\faIcon{check}] Écrivez lisiblement et justifiez vos réponses
			\item[\faIcon{check}] Respectez l'orthographe des nombres en lettres
			\item[\faIcon{check}] Les exercices peuvent être traités dans n'importe quel ordre
			\item[\faIcon{times}] Calculatrice non autorisée
		\end{itemize}
	\end{consignebox}
	
	\vspace{0.5cm}
	
	% Exercice 1
	\begin{exercisebox}{\faIcon{sort-numeric-down} Exercice 1 : Rang des chiffres \points{4}}
		
		Considérons le nombre : \textbf{7 256 394}
		
		\begin{enumerate}[label=\textcolor{primary}{\textbf{\alph*)}},leftmargin=*]
			\item Quel est le chiffre des centaines ? \dotfill \points{1}
			
			\item Quel est le chiffre des dizaines de milliers ? \dotfill \points{1}
			
			\item Quel est le nombre de milliers dans ce nombre ? \dotfill \points{1}
			
		%	\item Écris ce nombre dans le tableau de numération ci-dessous : \points{1}
%			\begin{center}
%				\renewcommand{\arraystretch}{2} % Augmente la hauteur des lignes
%				\begin{tabular}{|>{\centering\arraybackslash}p{1cm}|>{\centering\arraybackslash}p{1cm}|>{\centering\arraybackslash}p{1cm}||>{\centering\arraybackslash}p{1cm}|>{\centering\arraybackslash}p{1cm}|>{\centering\arraybackslash}p{1cm}||>{\centering\arraybackslash}p{1cm}|>{\centering\arraybackslash}p{1cm}|>{\centering\arraybackslash}p{1cm}|}
%					\hline
%					\multicolumn{3}{|c||}{Millions} & \multicolumn{3}{c||}{Milliers} & \multicolumn{3}{c|}{Unités} \\
%					\hline
%					c & d & u & c & d & u & c & d & u \\
%					\hline
%					\rule{0pt}{1.5em} & & & & & & & & \\ % Espace vertical pour écrire
%					\hline
%				\end{tabular}
%				\renewcommand{\arraystretch}{1} % Remet la valeur par défaut
%			\end{center}
			
			\item Donne la décomposition décimale de ce nombre : \points{1}
			
			\zonereponseLignes[0.5cm]{2}
		\end{enumerate}
	\end{exercisebox}
	
	% Exercice 2
	\begin{exercisebox}{\faIcon{spell-check} Exercice 2 : Écriture en lettres \points{5}}
		
		\textbf{Écris les nombres suivants en toutes lettres :}
		
		\begin{enumerate}[label=\textcolor{primary}{\textbf{\alph*)}},leftmargin=*]
			\item 384 : \dotfill \points{1}
			
			\item 1 200 : \dotfill \points{1}
			
			\item 5 087 : \dotfill 
			
			\dotfill \points{1}
			
			\item 14 680 : \dotfill 
			
			\dotfill \points{1}
			
			\item 300 092 : \dotfill 
			
			\dotfill \points{1}
		\end{enumerate}
	\end{exercisebox}
	
	% Exercice 3
	\begin{exercisebox}{\faIcon{ruler-horizontal} Exercice 3 : Demi-droite graduée \points{4}}
		
		\vspace{0.3cm}
		\textbf{Observe la demi-droite graduée ci-dessous :}
		
		\begin{center}
			\begin{tikzpicture}[scale=1]
				% Demi-droite
				\draw[thick,->] (0,0) -- (12,0);
				% Graduations principales
				\foreach \x in {0,2,4,6,8,10} {
					\draw[thick] (\x,0.1) -- (\x,-0.1);
					\node[below] at (\x,-0.3) {\ifnum\x=0 0\else \x0\fi};
				}
				% Graduations secondaires
				\foreach \x in {1,3,5,7,9,11} {
					\draw (\x,0.05) -- (\x,-0.05);
				}
				% Points
				\node[circle,fill=primary,inner sep=2pt,label=above:$A$] at (3,0) {};
				\node[circle,fill=primary,inner sep=2pt,label=above:$B$] at (7,0) {};
				\node[circle,fill=primary,inner sep=2pt,label=above:$C$] at (9.5,0) {};
			\end{tikzpicture}
		\end{center}
		
		\begin{enumerate}[label=\textcolor{primary}{\textbf{\alph*)}},leftmargin=*]
			\item Quelle est l'abscisse du point $A$ ? \dotfill \points{1}
			
			\item Quelle est l'abscisse du point $B$ ? \dotfill \points{1}
			
			\item Quelle est l'abscisse du point $C$ ? \dotfill \points{1}
			
			\item Place le point $D$ d'abscisse 45 sur la demi-droite ci-dessus. \points{1}
		\end{enumerate}
	\end{exercisebox}
	
	% Exercice 4
	\begin{exercisebox}{\faIcon{greater-than-equal} Exercice 4 : Comparaison et rangement \points{4}}
		
		\begin{enumerate}[label=\textcolor{primary}{\textbf{\alph*)}},leftmargin=*]
			\item Compare les nombres suivants en utilisant $<$, $>$ ou $=$ : \points{2}
			
			\begin{multicols}{4}
				\begin{itemize}[label=$\bullet$]
					\item $3\,456$ \dotfill $3\,546$
					\item $12\,089$ \dotfill $12\,098$
					\item $7\,000$ \dotfill $6\,999$
					\item $45\,230$ \dotfill $45\,203$
				\end{itemize}
			\end{multicols}
			
			\item Range les nombres suivants dans l'ordre croissant : \points{2}
			
			\begin{center}
				$8\,743$ \quad ; \quad $8\,347$ \quad ; \quad $8\,734$ \quad ; \quad $8\,437$ \quad ; \quad $8\,374$
			\end{center}
			
			\zonereponseLignes[0.5cm]{2}
		\end{enumerate}
	\end{exercisebox}
	
    %\newpage
    
	% Exercice 5
	\begin{exercisebox}{\faIcon{lightbulb} Exercice 5 : Problème \points{3}}
		
		\begin{tcolorbox}[colback=blue!5,colframe=primary!30,arc=2mm,boxrule=0.5pt]
			\textbf{Contexte :} La bibliothèque du collège compte ses livres.
			\begin{itemize}[noitemsep,topsep=5pt]
				\item Il y a 3 milliers de romans
				\item Il y a 15 centaines de livres documentaires
				\item Il y a 85 dizaines de bandes dessinées
			\end{itemize}
		\end{tcolorbox}
		
		\begin{enumerate}[leftmargin=*]
			\item \textbf{Calcule} le nombre total de livres dans la bibliothèque : \points{2}
			
			\zonereponseLignes[0.5cm]{3}
			
			\item \textbf{Écris} ce nombre total en toutes lettres : \points{1}
			
			\zonereponseLignes[0.5cm]{2}
		\end{enumerate}
	\end{exercisebox}
	
	\vspace{0.3cm}
	
	% Bonus
	\begin{tcolorbox}[
		colback=accent!10,
		colframe=accent,
		fonttitle=\bfseries,
		title={\faIcon{star} BONUS : Défi \badge{accent}{+2pts}},
		arc=3mm,
		boxrule=1.5pt,
		enhanced,
		width=\linewidth,
		grow to left by=3mm,
		grow to right by=3mm
		]
		Un nombre mystère a les propriétés suivantes :
		\begin{itemize}[leftmargin=*,noitemsep]
			\item C'est un nombre à 4 chiffres
			\item Le chiffre des milliers est 5
			\item Le chiffre des centaines est le double du chiffre des unités
			\item Le chiffre des dizaines est 3
			\item La somme de tous ses chiffres est 16
		\end{itemize}
		
		Trouve ce nombre mystère et écris-le en chiffres puis en lettres :
		
		\zonereponseLignes[0.5cm]{3}
	\end{tcolorbox}
	
	% Pied de page avec évaluation
	\vspace{0.3cm}
	\begin{center}
		\begin{tikzpicture}
			\node[draw=secondary,rounded corners=3mm,inner sep=8pt,fill=light!30] {
				\begin{tabular}{lccccc}
					\textbf{Auto-évaluation :} & 
					\faIcon{frown} Difficile & 
					\faIcon{meh} Moyen & 
					\faIcon{smile} Facile & 
					\phantom{xx} & 
					\textbf{Temps utilisé :} \underline{\phantom{xxxxx}} min
				\end{tabular}
			};
		\end{tikzpicture}
	\end{center}
	
	\newpage
	
	% =======================
	% CORRIGÉ
	% =======================
	
	\begin{center}
		\Large\textbf{CORRIGÉ}
	\end{center}
	
	\begin{exercisebox}{\faIcon{sort-numeric-down} Exercice 1 – Corrigé}
		Pour le nombre 7 256 394 :
		
		a) Chiffre des centaines : \rep{3}
		
		b) Chiffre des dizaines de milliers : \rep{5}
		
		c) Nombre de milliers : \rep{7 256} (car 7 256 394 = 7 256 $\times$ 1000 + 394)
		
		d) Tableau : 7 dans millions-unités, 2 dans milliers-centaines, 5 dans milliers-dizaines, 6 dans milliers-unités, 3 dans unités-centaines, 9 dans unités-dizaines, 4 dans unités-unités
		
		e) Décomposition : \rep{7 $\times$ 1~000~000 + 2 $\times$ 100~000 + 5 $\times$ 10~ 000 + 6 $\times$ 1~000 + 3 $\times$ 100 + 9 $\times$ 10 + 4}
	\end{exercisebox}
	
	\begin{exercisebox}{\faIcon{spell-check} Exercice 2 – Corrigé}
		a) 384 : \rep{trois-cent-quatre-vingt-quatre}
		
		b) 1 200 : \rep{mille-deux-cents}
		
		c) 5 087 : \rep{cinq-mille-quatre-vingt-sept}
		
		d) 14 680 : \rep{quatorze-mille-six-cent-quatre-vingts}
		
		e) 300 092 : \rep{trois-cent-mille-quatre-vingt-douze}
	\end{exercisebox}
	
	\begin{exercisebox}{\faIcon{ruler-horizontal} Exercice 3 – Corrigé}
		a) Abscisse de A : \rep{30}
		
		b) Abscisse de B : \rep{70}
		
		c) Abscisse de C : \rep{95}
		
		d) Le point D doit être placé entre les graduations 40 et 50, à mi-chemin (à 4,5 cm de l'origine sur le dessin)
	\end{exercisebox}
	
	\begin{exercisebox}{\faIcon{greater-than-equal} Exercice 4 – Corrigé}
		a) Comparaisons :
		• 3 456 \rep{<} 3 546
		• 12 089 \rep{<} 12 098
		• 7 000 \rep{>} 6 999
		• 45 230 \rep{>} 45 203
		
		b) Ordre croissant : \rep{8 347 < 8 374 < 8 437 < 8 734 < 8 743}
	\end{exercisebox}
	
	\begin{exercisebox}{\faIcon{lightbulb} Exercice 5 – Corrigé}
		a) Calcul :
		• Romans : 3 milliers = 3 000 livres
		• Documentaires : 15 centaines = 1 500 livres
		• BD : 85 dizaines = 850 livres
		Total : 3 000 + 1 500 + 850 = \rep{5 350 livres}
		
		b) En lettres : \rep{cinq-mille-trois-cent-cinquante}
	\end{exercisebox}
	
	\begin{tcolorbox}[colback=accent!10,colframe=accent]
		\textbf{BONUS – Corrigé}
		
		Analyse : Le nombre est de la forme 5c3u (milliers-centaines-dizaines-unités)
		• Chiffre des milliers : 5
		• Chiffre des dizaines : 3  
		• Chiffre des centaines = 2 $\times$ chiffre des unités, donc c = 2u
		• Somme des chiffres : 5 + c + 3 + u = 16
		
		Résolution : 5 + 2u + 3 + u = 16
		8 + 3u = 16
		3u = 8 → pas de solution entière
		
		Reconsidérons : Si u = 2, alors c = 4
		Vérification : 5 + 4 + 3 + 2 = 14 (incorrect)
		
		Si u = 3, alors c = 6
		Vérification : 5 + 6 + 3 + 3 = 17 (incorrect)
		
		La bonne réponse : u = \rep{2} et c = \rep{6}
		Le nombre est : \rep{5632}
		Vérification : 5 + 6 + 3 + 2 = 16 ✓
		
		En lettres : \rep{cinq-mille-six-cent-trente-deux}
	\end{tcolorbox}
	
\end{document}